% programme_summary_sections9_14.tex
% Sections 9--14 of the programme summary paper
% To be merged with programme_summary_section1.tex
% after sections 5--8.
%
% Convention: this file is self-contained but uses the same
% preamble and macros as Section 1.  For compilation, splice
% into the main document before \end{document}.


% ====================================================================
\section{The standard landscape}
\label{sec:landscape}
% ====================================================================

The bar construction, the genus tower, the Swiss-cheese
structure, and the collision residue are not abstractions
awaiting examples: they were \emph{found} by computing
examples.  The standard landscape is the testing ground
where the theory proves itself.

% ====================================================================
\subsection{Census}
\label{ssec:census}
% ====================================================================

Seven families span the landscape.  Each is chirally
Koszul.  The invariants computed in the monograph are
recorded in Table~\ref{tab:census}.

\begin{table}[ht]
\centering
\caption{The standard landscape: seven families and their
  bar-complex invariants}
\label{tab:census}
\renewcommand{\arraystretch}{1.7}
{\small
\begin{tabular}{lcccccl}
\toprule
\textbf{Algebra $\cA$}
  & $\boldsymbol{c(\cA)}$
  & $\boldsymbol{\kappa(\cA)}$
  & $\boldsymbol{r_{\max}}$
  & \textbf{Class}
  & $\boldsymbol{r(z)}$
  & $\boldsymbol{\cA^!}$ \\
\midrule
$\cH_k$
  & $1$
  & $k$
  & $2$ & $\mathbf{G}$
  & $k/z$
  & $\Sym^{\mathrm{ch}}(V^*)$ \\[2pt]
$\widehat{\fg}_k$
  & $\dfrac{kd}{k+h^\vee}$
  & $\dfrac{d(k+h^\vee)}{2h^\vee}$
  & $3$ & $\mathbf{L}$
  & $\Omega/((k{+}h^\vee)z)$
  & $\widehat{\fg}_{-k-2h^\vee}$ \\[6pt]
$\beta\gamma_\lambda$
  & $1-3(2\lambda{-}1)^2$
  & $c/2$
  & $4$ & $\mathbf{C}$
  & const.
  & $bc_\lambda$ \\[2pt]
$\Vir_c$
  & $c$
  & $c/2$
  & $\infty$ & $\mathbf{M}$
  & $\frac{c/2}{z^3} + \frac{2T}{z}$
  & $\Vir_{26-c}$ \\[4pt]
$\cW^k_3(\mathfrak{sl}_3)$
  & $c(k)$
  & $5c/6$
  & $\infty$ & $\mathbf{M}$
  & $r_{\cW_3}(z)$
  & $\cW^{k'}_3$ \\[2pt]
$\cW^k_N(\mathfrak{sl}_N)$
  & $c(k)$
  & $c \cdot (H_N{-}1)$
  & $\infty$ & $\mathbf{M}$
  & $r_{\cW_N}(z)$
  & $\cW^{k'}_N$ \\[2pt]
$V_\Lambda$\,(lattice)
  & $\mathrm{rank}(\Lambda)$
  & $\mathrm{rank}(\Lambda)$
  & $2$ & $\mathbf{G}$
  & $\mathrm{rank}/z$
  & $V_{\Lambda^*}$ \\
\bottomrule
\end{tabular}
}
\end{table}

\noindent
Here $d = \dim\fg$, $h^\vee$ is the dual Coxeter number,
$t = k + h^\vee$, $k' = -k - 2h^\vee$,
$H_N = \sum_{j=1}^{N} 1/j$ is the $N$-th harmonic number,
and $\Omega = \sum_a J^a \otimes J_a$ is the Casimir tensor.

The table has two readings.  Read horizontally, each row is
a complete portrait: central charge, modular characteristic,
shadow depth, collision residue, Koszul dual.  Read
vertically, the columns expose the structural regularities
that no single example can reveal.  The modular characteristic
$\kappa$ is always a rational function of the level;
the complementarity sum $\kappa + \kappa^!$ vanishes for
free-field and Kac--Moody algebras, equals $13$ for
Virasoro, and is $(c + c')(H_N - 1)$ for
$\cW_N$.  The collision residue $r(z)$ is scalar for
classes~$\mathbf{G}$ and~$\mathbf{C}$, matrix-valued for
classes~$\mathbf{L}$ and~$\mathbf{M}$.

% ====================================================================
\subsection{The DS reduction chain}
\label{ssec:ds-chain}
% ====================================================================

Drinfeld--Sokolov reduction links the families.  For
$\fg = \mathfrak{sl}_N$, the principal DS reduction
\begin{equation}\label{eq:ds-chain}
V_k(\mathfrak{sl}_N)
\;\xrightarrow{\;\mathrm{DS}\;}
\;\cW^k_N(\mathfrak{sl}_N)
\end{equation}
transforms the shadow obstruction tower in a controlled
way:
\begin{enumerate}[(i)]
\item \emph{Depth increase.}
  The affine algebra $V_k(\mathfrak{sl}_N)$ is
  class~$\mathbf{L}$ ($r_{\max} = 3$); the
  $\cW$-algebra $\cW^k_N$ is class~$\mathbf{M}$
  ($r_{\max} = \infty$).  DS reduction creates the
  infinite shadow tower from finite data.

\item \emph{Quartic creation from zero.}
  The quartic shadow $S_4$ vanishes for all affine
  algebras (the Jacobi identity kills the
  obstruction); after DS reduction,
  $S_4(\cW^k_N) \ne 0$ for generic~$k$.
  For Virasoro:
  $S_4(\Vir_c) = 10/[c(5c+22)]$.

\item \emph{Central charge is additive; $\kappa$ is not.}
  The BRST ghost system contributes
  $c_{\mathrm{ghost}} = N(N{-}1)$ to the central
  charge but a $k$-dependent correction to~$\kappa$.
  The commutation failure of DS reduction with the
  shadow map is the source of all non-formality in
  the $\cW$-algebra landscape.
\end{enumerate}

The chain $\cH \to \widehat{\fg}_k \to \cW^k_N$ is a
cascade of increasing shadow complexity:
$\mathbf{G} \to \mathbf{L} \to \mathbf{M}$, with each
reduction introducing new structure that the previous
level could not see.


% ====================================================================
\subsection{The multi-weight correction}
\label{ssec:multi-weight-landscape}
% ====================================================================

For algebras with generators of distinct conformal
weights, the genus expansion
receives a correction.  The $\cW_3$ algebra has generators
$T$ (weight~$2$) and $W$ (weight~$3$); at genus~$2$ the
cross-channel graph sum contributes
\begin{equation}\label{eq:delta-f2-w3}
\delta F_2(\cW_3)
\;=\;
\frac{c + 204}{16c}
\;>\; 0
\qquad
\text{for all } c > 0\,.
\end{equation}
The correction arises from boundary graphs
of~$\overline{\cM}_{2,0}$ whose edges carry mixed
propagator channels (weight-$2$ and weight-$3$ generators
sewed through different Hodge data at the vertex level).
The propagator itself is always weight-$1$ at the edge
level ($d\log E$ has conformal weight~$1$); the mixing
enters through the vertex structure constants.

For uniform-weight algebras (all generators of the same
conformal weight), $\delta F_g^{\mathrm{cross}} = 0$
at all genera and the scalar universality
$F_g = \kappa \cdot \lambda_g^{\mathrm{FP}}$ holds
without correction.  For multi-weight algebras, the
full genus expansion is
$F_g = \kappa \cdot \lambda_g^{\mathrm{FP}}
+ \delta F_g^{\mathrm{cross}}$,
with
$\delta F_g^{\mathrm{cross}}$ computable from the
mixed-channel graph sum.



% ====================================================================
\section{The MC programme}
\label{sec:mc-programme}
% ====================================================================

The MC programme converts the geometric structure of the
bar complex into five concrete mathematical problems, each
asking whether a structural property of the universal MC
element~$\Theta_\cA$ persists under a specific operation:
PBW concentration, bar-intrinsic existence, categorical
generation, inverse-limit completion, and analytic sewing.
All five are proved.

% ====================================================================
\subsection{MC1: PBW concentration}
\label{ssec:mc1}
% ====================================================================

The PBW spectral sequence of~$B(\cA)$ collapses at~$E_2$
for every freely strongly generated chiral algebra.  Free
strong generation implies the PBW filtration on the bar
complex is exhaustive, and the associated graded is a
cofree coalgebra whose cohomology concentrates in bar
degree~$1$.  This is chirally Koszul: the bar cohomology
$H^*(B(\cA))$ is the Koszul dual $\cA^!$ (as a graded
coalgebra).

The proof applies to the universal $\cW$-algebra
$\cW^k(\fg)$ at all levels~$k$ (Feigin--Frenkel free
generation).  The simple quotient $L_k(\fg)$ at admissible
levels has Koszulness proved for $\fg = \mathfrak{sl}_2$
at all admissible levels; higher rank is open.

% ====================================================================
\subsection{MC2: bar-intrinsic existence}
\label{ssec:mc2}
% ====================================================================

The universal MC element $\Theta_\cA = D_\cA - d_0$
exists for all modular Koszul chiral algebras, with no
restriction to simple Lie symmetry.  The proof is a
tautology at the right level of abstraction: $D_\cA^2 = 0$
implies $\Theta_\cA$ satisfies the MC
equation~(\ref{eq:mc-equation}).  The algebraic-family
rigidity theorem extends the scalar identification
$\Theta_\cA^{\min} = \eta \otimes \Gamma_\cA$ to all
algebraic families with rational OPE coefficients, covering
the entire standard Lie-theoretic landscape at all
non-critical levels.

% ====================================================================
\subsection{MC3: categorical generation}
\label{ssec:mc3}
% ====================================================================

For every simple Lie algebra~$\fg$, the Drinfeld--Kazhdan
category $\mathrm{DK}(\fg)$ is thickly generated by
evaluation modules.  The proof combines the chromatic
filtration on $\mathrm{Rep}(Y(\fg))$ with the categorical
Clebsch--Gordan closure via multiplicity-free
$\ell$-weights (Chari--Moura), the Francis--Gaitsgory
pro-nilpotent completion, and Drinfeld--Kazhdan
compactness.  The type~A proof uses the minuscule
hypothesis; the all-types generalization replaces it with
the strictly weaker multiplicity-free $\ell$-weight
property, which holds for all simple types.

% ====================================================================
\subsection{MC4: inverse-limit completion}
\label{ssec:mc4}
% ====================================================================

The bar-cobar adjunction extends to inverse limits via the
strong completion-tower theorem: the strong filtration axiom
$\mu_r(F^{i_1}, \ldots, F^{i_r})
\subset F^{i_1 + \cdots + i_r}$
forces an arity cutoff that makes continuity and the
Mittag-Leffler condition automatic.  The completion closure
$\mathrm{CompCl}(F_{\mathrm{ft}})$ carries a quasi-inverse
bar-cobar homotopy equivalence, stable under MC twisting.

The programme splits: MC4$^+$ (positive towers:
$\cW_{1+\infty}$, affine Yangians) is solved by weight
stabilization.  MC4$^0$ (resonant towers: Virasoro,
non-quadratic $\cW_N$) is reduced to a finite resonance
problem by the resonance-filtered bar-cobar theorem.

% ====================================================================
\subsection{MC5: analytic sewing}
\label{ssec:mc5}
% ====================================================================

The sewing envelope $\cA^{\mathrm{sew}}$, the Hausdorff
completion of~$\cA$ for the all-amplitude seminorm
topology, carries the analytic data needed for genus-$g$
partition functions.  The HS-sewing criterion provides a
sufficient condition: polynomial OPE growth and
subexponential sector growth imply convergence at all
genera.  Every standard family satisfies this condition.

For Heisenberg, the sewing theorem is explicit: the
one-particle Bergman reduction writes the genus-$g$
partition function as a Fredholm determinant.  For general
modular Koszul algebras, the HS-sewing criterion gives
convergence without an explicit closed form.  The BV/BRST
identification of sewing amplitudes with bar complex
data at higher genus remains conjectural.


% ====================================================================
\section{Arithmetic}
\label{sec:arithmetic}
% ====================================================================

The shadow obstruction tower produces arithmetic invariants
from vertex-algebraic data.  These invariants connect the
bar complex to number theory through three structures: a
Dirichlet series, a categorical zeta function, and a depth
decomposition.

% ====================================================================
\subsection{The shadow Eisenstein theorem}
\label{ssec:eisenstein}
% ====================================================================

The genus-$1$ amplitude of a modular Koszul chiral
algebra~$\cA$ is the quasi-modular Eisenstein series
$\kappa(\cA) \cdot E_2^*(\tau)$, where
$E_2^*(\tau) = E_2(\tau) - 3/(\pi \operatorname{Im}\tau)$
is the completed weight-$2$ Eisenstein series.  The
Fourier coefficients $\sigma_1(n)$ of $E_2^*$ produce
a Dirichlet series:
\begin{equation}\label{eq:dirichlet-series}
D_2(\cA, s)
\;:=\;
-24\kappa(\cA) \cdot
\zeta(s) \cdot \zeta(s - 1)\,,
\end{equation}
an Eisenstein $L$-function with poles at $s = 1$ (from
$\zeta(s)$) and $s = 2$ (from $\zeta(s-1)$).  The identity
follows from the $\sigma_1$ Dirichlet convolution applied to
the Fourier coefficients of $\kappa \cdot E_2^*(\tau)$.

The genus-$1$ propagator is $E_2^*$, not a holomorphic
modular form: $E_2^*$ is quasi-modular of weight~$2$ and
depth~$1$, transforming with an additive anomaly
$3\tau/(\pi i)$ under $\mathrm{SL}(2,\mathbb{Z})$.
Products of $E_2^*$ live in the ring of quasi-modular
forms $\mathbb{C}[E_2^*, E_4, E_6]$, not in the ring of
holomorphic modular forms $\mathbb{C}[E_4, E_6]$.  This
distinction is load-bearing: holomorphic dimension formulas
(e.g., $\dim S_k = 0$ for $k < 12$) do not apply to
quasi-modular objects.


% ====================================================================
\subsection{The categorical zeta function}
\label{ssec:categorical-zeta}
% ====================================================================

The dimension-counting zeta function of
$\mathfrak{sl}_2$-representations is
\begin{equation}\label{eq:cat-zeta}
\zeta^{\mathrm{DK}}_{\mathfrak{sl}_2}(s)
\;:=\;
\sum_{n \ge 1} (\dim V_n)^{-s}
\;=\;
\sum_{n \ge 1}(n+1)^{-s}
\;=\;
\zeta(s) - 1\,.
\end{equation}
The Riemann zeta function (shifted by~$1$) counts
$\mathfrak{sl}_2$-representations weighted by inverse
dimension.  The link to the shadow--DK bridge is that
the categorical zeta encodes the multiplicity structure
of the evaluation modules whose thick generation proves
MC3.

For $\mathfrak{sl}_N$ with $N \ge 3$,
$\zeta^{\mathrm{DK}}_{\mathfrak{sl}_N}(s) :=
\sum_\lambda (\dim V_\lambda)^{-s}$ is not
multiplicative: the Weyl dimension formula is a product
of linear forms, but the sum over the dominant chamber
has no Euler product.  The $N = 2$ coincidence with
$\zeta(s)$ is a low-rank accident.

% ====================================================================
\subsection{Depth decomposition}
\label{ssec:depth}
% ====================================================================

The total shadow depth decomposes:
\begin{equation}\label{eq:depth-decomposition}
d(\cA)
\;=\;
1 + d_{\mathrm{arith}}(\cA) + d_{\mathrm{alg}}(\cA)\,,
\end{equation}
where $d_{\mathrm{arith}}$ counts arithmetic contributions
(cusp forms in the automorphic decomposition of the shadow
generating function) and $d_{\mathrm{alg}}$ counts
algebraic contributions (non-formality of the $L_\infty$
structure).  The algebraic depth
$d_{\mathrm{alg}} \in \{0, 1, 2, \infty\}$ recovers
the G/L/C/M classification: $d_{\mathrm{alg}} = 0$
for class~$\mathbf{G}$, $d_{\mathrm{alg}} = 1$ for
class~$\mathbf{L}$, $d_{\mathrm{alg}} = 2$ for
class~$\mathbf{C}$,
$d_{\mathrm{alg}} = \infty$ for class~$\mathbf{M}$.
Shadow depths $\ge 5$ are purely arithmetic: they arise
from cusp forms, beginning with the weight-$12$ cusp
form $\Delta_{12}$ at the genus-$2$ level.

The decomposition separates the algebraic content
(governed by the MC equation on the cyclic deformation
complex) from the arithmetic content (governed by
automorphic forms on moduli of curves).  The two are
independent: a class-$\mathbf{G}$ algebra (algebraically
trivial) can have nontrivial arithmetic depth, and a
class-$\mathbf{M}$ algebra (algebraically infinite)
can have vanishing arithmetic depth at low genus.


% ====================================================================
\section{Frontier directions}
\label{sec:frontier}
% ====================================================================

Three leaps organize the frontier.  The first extends the
modular Koszul programme from chiral algebras to
holographic systems.  The second extends it from formal
algebraic data to analytic sewing amplitudes.  The third
extends it from $2$-dimensional geometry to
$3$-dimensional topology.

% ====================================================================
\subsection{First leap: holographic modular Koszul data}
\label{ssec:holographic}
% ====================================================================

Every holomorphic-topological system~$T$ in the sense of
Costello--Li should be controlled by a \emph{holographic
modular Koszul datum}
\begin{equation}\label{eq:holographic-datum}
\cH(T)
\;=\;
(\cA,\, \cA^!,\, \mathfrak{C},\, r(z),\,
\Theta_\cA,\, \nabla^{\mathrm{hol}})\,,
\end{equation}
consisting of the boundary algebra, its Koszul dual, the
bulk algebra (chiral derived centre), the collision
residue, the universal MC element, and the holographic
shadow connection.  The collision residue is the genus-$0$
binary shadow of~$\Theta_\cA$:
$r(z) = \Res^{\mathrm{coll}}_{0,2}(\Theta_\cA)$.

This is the Platonic holographic programme.  Its five
theorem targets are: boundary-defect realization (the
boundary algebra is the algebra of boundary conditions
of the $3$d theory), Yangian-shadow identification (the
Yangian controls line operators), sphere reconstruction
(genus-$0$ amplitudes recover the boundary OPE), quartic
resonance obstruction (the quartic shadow controls the
deformation quantization), and singular-fibre descent (the
critical-level limit produces the Feigin--Frenkel centre).

Two other directions in the first leap: the
factorization-envelope technology of Nishinaka and Vicedo,
which gives a functorial pipeline from Lie conformal
algebras to vertex algebras; and the non-principal
$\cW$-algebra corridor, where hook-type nilpotent orbits
in type~A provide the first proved non-principal Koszul
dualities.


% ====================================================================
\subsection{Second leap: analytic sewing}
\label{ssec:sewing-frontier}
% ====================================================================

The algebraic bar construction produces formal power
series in the sewing parameters.  The analytic sewing
programme asks: when do these series converge?

The answer is the sewing envelope $\cA^{\mathrm{sew}}$,
the Hausdorff completion for the all-amplitude seminorm
topology.  The HS-sewing condition (polynomial OPE growth
and subexponential sector growth) is satisfied by the
entire standard landscape.  The Heisenberg sewing theorem
gives an explicit Fredholm determinant.  The lattice
sewing envelope is constructed; the analytic realization
criterion (a sufficient condition for passage from
formal to analytic) is conjectural.

At genus $g \ge 1$ the curvature~\eqref{eq:curvature}
forces passage from ordinary to coderived categories:
the bar complex is no longer flat, and the correct
homological framework is the coderived category of
curved coalgebras.  The analytic programme at higher genus
requires the coderived analytic shadow
$Q_g^{\mathrm{an}}(\cA)$.


% ====================================================================
\subsection{Third leap: three-dimensional topology}
\label{ssec:3d}
% ====================================================================

The Swiss-cheese operad $\mathrm{SC}^{\mathrm{ch,top}}$
packages both the chiral (holomorphic, $\mathbb{C}$)
and topological ($\mathbb{R}$) directions of a
$3$-dimensional holomorphic-topological theory.  The
bar complex factorizes along both: the differential
encodes the $\mathbb{C}$-direction factorization, the
coproduct encodes the $\mathbb{R}$-direction
factorization.  The ordered bar
$\barB^{\mathrm{ord}}(\cA)$ is simultaneously an
$E_1$-coalgebra (from the linear order on~$\mathbb{R}$)
and an $E_\infty$-coalgebra (from the factorization
structure on~$\mathbb{C}$).

The third leap is the passage from the algebraic
Swiss-cheese structure (proved) to the geometric
realization on $3$-manifolds of the form
$\Sigma_g \times [0,1]$.  The $\mathbb{C}$-direction
carries the chiral algebra; the $\mathbb{R}$-direction
carries the Yangian / quantum group.  The annulus trace
$\Delta_{\mathrm{ns}}(\mathrm{Tr}_\cA) =
\kappa \cdot \lambda_1$
provides the first open-to-closed map at genus~$1$;
the full genus tower progressively restores
bidirectionality between open and closed sectors.

The deepest open problem in this direction is CY-A at
$d = 3$ (Section~\ref{sec:open-problem}).


% ====================================================================
\section{The single open problem}
\label{sec:open-problem}
% ====================================================================

Everything in Sections~\ref{sec:bar}--\ref{sec:frontier}
rests on two-dimensional geometry: curves, their moduli
spaces, and the factorization structure of colliding points
on a curve.  The three-dimensional extension is conditional
on a single unresolved input.

\begin{conjecture}[CY-A at $d = 3$]\label{conj:cy-a}
The Calabi--Yau/$A_\infty$ correspondence extends from
$d = 2$ (proved) to $d = 3$: the chain-level
$S^3$-framing of the topological direction in a
$3$-dimensional holomorphic-topological theory produces
an $A_\infty$-structure on the boundary algebra that is
compatible with the chiral bar complex.
\end{conjecture}

At $d = 2$, the CY-A correspondence is a theorem: the
Fukaya category of a symplectic surface carries an
$A_\infty$-structure whose bar construction is the wrapped
Floer complex, and the $A_\infty$-formality on the Koszul
locus is the $E_2$-collapse of the bar spectral sequence.
At $d = 3$, the analogue requires the chain-level framing
of the $S^3$-fibre in the conifold geometry to be
compatible with the factorization structure on the
$\mathbb{C}$-direction.  This is a statement about the
interaction of holomorphic and topological data at the
chain level, not merely at the level of cohomology.

The difficulty is not conceptual but technical: the
$d = 2$ proof uses the Riemann surface structure in an
essential way (the Abel--Jacobi map, the period matrix,
the prime form), and the $d = 3$ analogue must replace
surface geometry with $3$-manifold geometry while
preserving the $A_\infty$-coherences.

No other open conjecture in the programme has unresolved
algebraic prerequisites.  The five MC problems are solved;
the Koszulness characterization is complete to within
the D-module purity converse; the shadow obstruction tower
is constructed at all arities; the arithmetic invariants
are computed; the standard landscape is fully mapped.
CY-A at $d = 3$ is the single point where the programme
reaches beyond what it can currently prove, and it is the
natural boundary: the passage from the algebraic curve
to the three-manifold.


% ====================================================================
\section*{References}
\label{sec:references}
% ====================================================================

\begin{thebibliography}{99}

\bibitem{Arnold69}
V.\,I.~Arnold,
The cohomology ring of the group of dyed braids,
\textit{Mat.\ Zametki}~\textbf{5} (1969), 227--231.

\bibitem{BD04}
A.~Beilinson and V.~Drinfeld,
\textit{Chiral Algebras},
AMS Colloquium Publications, vol.~51,
American Mathematical Society, Providence, RI, 2004.

\bibitem{BGS96}
A.~Beilinson, V.~Ginzburg, and W.~Soergel,
Koszul duality patterns in representation theory,
\textit{J.~Amer.\ Math.\ Soc.}~\textbf{9} (1996),
473--527.

\bibitem{CG17}
K.~Costello and O.~Gwilliam,
\textit{Factorization Algebras in Quantum Field Theory},
vols.~1--2, Cambridge University Press, 2017/2021.

\bibitem{CL16}
K.~Costello and S.~Li,
Twisted supergravity and its quantization,
preprint, \texttt{arXiv:1606.00365}, 2016.

\bibitem{CM08}
V.~Chari and A.~Moura,
Characters and blocks for finite-dimensional
representations of quantum affine algebras,
\textit{Int.\ Math.\ Res.\ Not.}\ (2004), no.~5,
257--298.

\bibitem{DK84}
V.\,G.~Drinfeld and V.\,G.~Kazhdan,
Algebraic geometry of the space of quasi-maps
and the structure of representations of quantum groups,
unpublished notes, 1984.

\bibitem{Dri85}
V.\,G.~Drinfeld,
Hopf algebras and the quantum Yang--Baxter equation,
\textit{Dokl.\ Akad.\ Nauk SSSR}~\textbf{283} (1985),
1060--1064.

\bibitem{FBZ04}
E.~Frenkel and D.~Ben-Zvi,
\textit{Vertex Algebras and Algebraic Curves},
2nd ed., Mathematical Surveys and Monographs, vol.~88,
American Mathematical Society, Providence, RI, 2004.

\bibitem{FF92}
B.~Feigin and E.~Frenkel,
Affine Kac--Moody algebras at the critical level and
Gel'fand--Diki\u{\i} algebras,
\textit{Int.\ J.\ Mod.\ Phys.\ A}~\textbf{7}
(suppl.\ 1A) (1992), 197--215.

\bibitem{FG12}
J.~Francis and D.~Gaitsgory,
Chiral Koszul duality,
\textit{Selecta Math.\ (N.S.)}~\textbf{18} (2012),
27--87.

\bibitem{FM94}
W.~Fulton and R.~MacPherson,
A compactification of configuration spaces,
\textit{Ann.\ of Math.\ (2)}~\textbf{139} (1994),
183--225.

\bibitem{FP00}
C.~Faber and R.~Pandharipande,
Hodge integrals and moduli spaces of curves,
\textit{Enseign.\ Math.\ (2)}~\textbf{48} (2002),
173--195.

\bibitem{Gai16}
D.~Gaiotto,
Twisted compactifications of 3d
$\mathcal{N}=4$ theories and conformal blocks,
\textit{J.\ High Energy Phys.}\ (2019), no.~2, 61.

\bibitem{GL19}
D.~Gaiotto and J.~Lamy-Poirier,
Holomorphic-topological twists,
unpublished notes, 2019.

\bibitem{KRW03}
V.~Kac, S.-S.~Roan, and M.~Wakimoto,
Quantum reduction for affine superalgebras,
\textit{Comm.\ Math.\ Phys.}~\textbf{241} (2003),
307--342.

\bibitem{LV12}
J.-L.~Loday and B.~Vallette,
\textit{Algebraic Operads},
Grundlehren der Mathematischen Wissenschaften,
vol.~346, Springer, Berlin, 2012.

\bibitem{Mok25}
L.~Mok,
Log Fulton--MacPherson compactifications and
tropicalization of configuration spaces,
preprint, 2025.

\bibitem{Mor26}
Y.~Moriwaki,
Conformally flat factorization homology in
$\mathrm{IndHilb}$,
preprint, 2026.

\bibitem{MS24}
N.~Markarian and H.~L.~Tanaka,
Homotopy chiral algebras,
preprint, 2024.

\bibitem{Nis25}
D.~Nishinaka,
Factorization envelopes of Lie conformal algebras,
preprint, 2025.

\bibitem{Priddy70}
S.~Priddy,
Koszul resolutions,
\textit{Trans.\ Amer.\ Math.\ Soc.}~\textbf{152}
(1970), 39--60.

\bibitem{PTVV13}
T.~Pantev, B.~To\"en, M.~Vaqui\'e, and G.~Vezzosi,
Shifted symplectic structures,
\textit{Publ.\ Math.\ Inst.\ Hautes \'Etudes Sci.}
\textbf{117} (2013), 271--328.

\bibitem{RNW19}
D.~Robert-Nicoud and F.~Wierstra,
Homotopy morphisms between convolution homotopy
Lie algebras,
\textit{J.\ Noncommut.\ Geom.}~\textbf{13} (2019),
1463--1520.

\bibitem{Val16}
B.~Vallette,
Homotopy theory of homotopy algebras,
\textit{Ann.\ Inst.\ Fourier}~\textbf{70} (2020),
683--738.

\bibitem{Vic25}
B.~Vicedo,
Non-chiral factorization algebras and vertex algebras,
preprint, 2025.

\bibitem{WK84}
M.~Wakimoto and V.~Kac,
Quantum reduction and representation theory
of superconformal algebras,
\textit{Adv.\ Math.}~\textbf{185} (2004), 400--458.

\end{thebibliography}

\bigskip
\noindent\hrulefill

\medskip
\noindent
Modular Koszul duality for chiral algebras is $E_1$--$E_1$
operadic Koszul duality transported into the homotopical
modular chiral realm on algebraic curves.

\end{document}
